\chapter{序論}
\label{chap:序論}

\section{背景と目的}
\label{sec:背景と目的}

\subsection{ロボットナビゲーションの進化と課題}
近年，自律移動ロボットのナビゲーション技術は Vision-Language-Action（VLA）に代表される枠組みにより急速に発展している．
VLA は視覚情報と言語指示を統合し，「キッチンへ行って」「赤いボトルのある場所まで移動して」といった自然言語による指示から自律的に行動を生成する．
また，VLA の具体例として NaVILA \cite{NaVILA_paper} が提案されている．

しかし，既存の VLA モデルの多くは目的地への到達に主眼を置いており，
「どのように移動するか」という\textbf{移動のスタイル}の制御は十分に扱われていない．
実環境では効率性だけでなく，静粛性や慎重さ，力強さといった運動の質も重要である．
病院内では静かに慎重に振る舞い，急ぐ場面ではきびきびと移動するといった「振る舞いの質感」は，
従来の「Go to [Place]」型の指示体系では表現しにくかった．
こうした「移動の仕方」を専門的なパラメータ調整なしに直感的に伝えられるインターフェースが必要である．
例えば，病院・介護施設や図書館のように\textbf{静粛性}が求められる場面では「そろりそろり」といった表現が有効である．
一方で緊急時の迅速移動では，「さっさ」といった\textbf{素早い移動}を指示する語が直感的である．
このようにオノマトペは，場所・状況・人との距離感に応じた「望ましい移動の質感」を短い語で示せるため，非専門家でもその場にふさわしい振る舞いを簡潔に指示できる点で有用である．

\subsection{オノマトペによる直感的なスタイル記述}
本研究では\textbf{オノマトペ（擬音語・擬態語）}を用いて移動スタイルを記述・制御する手法を提案する．
日本語のオノマトペは「てくてく」「のっしのっし」「そろりそろり」など短い語句で，
速度やリズム，力強さ，滑らかさといった多次元的な運動特徴を表現できる（音象徴性 \cite{hoyo}）．
この性質を活かせば，ユーザは数値パラメータを指定せずとも感性的な語彙で運動の質を伝えられる．
VLA が担う空間的到達（Where）に対し，オノマトペは移動の仕方（How）を補完する言語表現として有効である．

\subsection{共有潜在空間による言語と動作の統合}
オノマトペに応じたスタイル動作をロボットに生成させるには，言語と身体動作を対応付ける表現が必要となる．
従来はモーションキャプチャ（MoCap）による模倣学習が主流であったが，
オノマトペが内包する微妙なニュアンスを網羅するデータ収集にはコストがかかる．

そこで本研究では，\textbf{動画データ}から 2D キーポイントを推定して得られるオノマトペ付きのモーションデータを活用し，
言語と動作を結ぶ\textbf{共有潜在空間}を構築する．
具体的には MotionCLIP \cite{MotionCLIP} を基盤とし，
オノマトペに対応する二次元の歩行モーションデータ（HOYO データセット \cite{hoyo}）とオノマトペのテキスト埋め込みを同一の潜在空間へ写像する．
また学習データにはない未知語への対応を考慮し，まず言語と動作の対応を対照学習する．

さらに，この潜在空間上の類似度を報酬として強化学習を行い，
物理的に妥当な歩行を維持しつつオノマトペのニュアンスを反映した動作を獲得する．
これにより「到達（Where）」と「移動様式（How）」を統合した直感的なナビゲーション指示を目指す．

\subsection{本研究の目的と新規性}
本研究の目的は，オノマトペで移動スタイルを指定できるナビゲーションの実現である．
ユーザの感性に沿った「振る舞いの質」をロボットの行動生成に反映する．
新規性は，従来の VLA が主に扱ってきた目的地到達に加え，
\textbf{日本語オノマトペという直感的な言語表現で移動様式を制御し}，
対照学習と強化学習を組み合わせて実現する点にある．

\section{本研究のアプローチ}
\label{sec:approach}

本研究では以下の手順でオノマトペに基づくスタイル付与ナビゲーションを実現する．
\begin{enumerate}
  \item \textbf{データ準備と埋め込みの選定}：
        HOYO の 2D 歩容系列を前処理し，オノマトペのテキスト埋め込み（Sarashina / SigLIP）を用意する．
  \item \textbf{共有潜在空間の構築（Phase 1）}：
        MotionCLIP を基盤に VAE による再構成と教師あり対照学習で言語と動作の対応を学習する．
  \item \textbf{強化学習によるスタイル制御（Phase 2）}：
        潜在空間上の類似度を報酬とし，PPO でオノマトペに合ったスタイル歩行を獲得する．
  \item \textbf{実験設計と評価}：
        埋め込みモデルの比較，対照学習の検索指標（R@K）/MPJPE/シルエット係数，強化学習の前進維持（補助速度条件）・転倒率・スタイル一致度・語彙間の相対速度差で評価する．
\end{enumerate}

\section{本論文の構成}
本論文は以下のように構成される．

第\ref{chap:theory_related_work}章ではナビゲーション技術とオノマトペ・音象徴性に関する研究を概観し，本研究の位置付けを示す．

第\ref{chap:proposed_method}章では提案手法を述べる．HOYO の前処理，Phase 1 の対照学習，Phase 2 の強化学習の順に説明する．

第\ref{chap:実験}章では埋め込み比較，対照学習，強化学習の実験設定と評価指標を整理する．

第\ref{chap:結果}章では対照学習と強化学習の結果を定量・定性の両面から示す．

第\ref{chap:考察}章ではオノマトペによるスタイル制御の特性と限界を議論する．

第\ref{chap:結論}章で本研究を総括し，今後の展望を述べる．
